\section{考察と限界}\label{sec:discussion}

\subsection{固定憲法と適応する制度}\label{sec:discussion:p2}

本方式は、無制限の自己改変を意図的に採用しない。憲法、遷移規則、
権限表、固定検証器は実行中に進化しない。そのため、まったく新しい役割
や遷移を導入するには人間の改正が必要となる。一方、この制限があるから
こそ制度変更を解釈できる。候補は既知の権限範囲内で改善できるが、自分
を裁く範囲そのものを広げることはできない。

憲法自体が誤っている可能性は残る。規則のハッシュ値を固定すれば変更は
明示できるが、規則が賢明であることまでは保証しない。したがって、人間
による査読、改正理由、回帰試験の証拠は、信頼すべき手続きの一部であり
続ける。Constitutional ARI-RQGM が減らすのはモデル出力へ与える権限
であり、法を定める人間の責任ではない。

緊急隔離は期内固定を維持する。T16 は現期を強制終了し、違反要素を隔離
しながら新しい指紋値の期を不可分に開始する。これにより、再生時の封じ込め
位置が明確になり、一つの指紋値が二つの稼働集合を表すことを防ぐ。ただし、
裁定経路への妨害や、緊急取引内での代替役確保は未解決である。

\subsection{権限境界での指示追従}\label{sec:discussion:seams}

危険な故障の多くは、モデルが生成した記録を決定論的な状態へ反映する
境界で起きる。形式検査だけでは足りない。形式が正しくても、誤った役割
が作成し、不正な遷移を要求し、権限を広げ、封印された文章を参照し、
修復に必要な来歴を欠くことがある。カーネルは、構造、作成者、時期、
権限、ハッシュ値、閲覧範囲を別々に検査する。

意味内容に関する限界もある。形式上有効な攻撃や裁定であっても、結論が
誤っている可能性はある。役割分離と過去事例による再試験は、誤りを観測
して異議を申し立てやすくするが、モデルの判断を真実へ変えるわけでは
ない。実行証拠や人手で整備した正解がある箇所では固定基準を使い、それ
以外の結論は不確実な判断として扱う必要がある。

同様に、「統治に失敗しても研究を止めない」は一般的な安全原則ではなく、
可用性を選ぶ条件付き方針である。継続してよいのは、隔離環境内の候補
生成と局所計算に限る。外部サービスへの書込み、実験装置の操作、本番
データの変更など不可逆な行為には、統治系と独立した固定許可口を置き、
統治故障時は閉じる必要がある。この条件は現カーネルが強制していない。

\subsection{規模拡大の限界}\label{sec:discussion:scaling}

統治処理は、モデル呼び出し、過去事例の再試験、記憶容量、期の境界での
待ち時間を増やす。予算が不足した場合は統治の強度を下げ、任意の共進化
を抑えることができるが、制度を十分に調べない代償が生じる。また、台帳
には候補、役割間の関係、修復対象となる依存関係が蓄積するため、長期
運用では索引、圧縮、保存方針が必要になる。

第二版の監査ハッシュ値は、種類、識別子、取引、内容、直前値を結び、
再生の意味に関わる途中改変を示せるが、履歴と局所固定値を同じ
過去の接頭辞へ戻す攻撃は示せない。強い巻戻し耐性には、外部に保存した
署名付き先頭値、独立監査者への定期固定、または単調計数器が必要である。
また、永久の追記保存は機密情報、個人情報、削除要求と緊張する。内容を
暗号化して鍵を破棄する方式や、監査可能な秘匿化記録を設計しない限り、
機微な原文を台帳へ格納すべきではない。

現実装は、一つの保存点につき正本となる出来事順序が一つだけ存在すると
仮定する。分散した研究実行系へ広げるには、複製された記録の整合、
競合の決定論的解決、同時生成された成果物へどの期の制度を適用するか
という定義が要る。現行の単一書き手は局所的には健全だが、分散合意の
手順ではない。

比較用の資料も制約となる。査読役の改善には未使用の基準資料、隔離再
生成には十分に抽象化した失敗要約、効用規則の比較には安定した評価軸が
必要である。これらが欠ける場合、仕組みは優越性を作り上げず、現職維持
へ縮退する。この方針は完全性を守る一方、適応の機会を減らす。

同じ検証済み攻撃集合で候補と効用規則を反復選択すれば、既知事例へ
過適合しうる。信頼できる評価には、開発用、選択用、秘匿した最終事例を
分け、適応的な反復選択を考慮し、同一モデル・異モデル・異提供者の役割
配置を比較する必要がある。記録上の役割分離は、同じモデル系列を使う
役割間の誤りを独立にはしない。結託、裁定妨害、偏った基準資料は未解決
の実証課題である。

\subsection{今後の課題}\label{sec:discussion:outlook}

次の実証段階では、統治設定、課題群、モデル、乱数種を変えた比較を
事前登録して行う必要がある。次の実装段階では、証拠の被覆範囲を明示し、
各役割に実際に適用された制裁、再試験事例、基準資料、後継候補の段階を
報告すべきである。効用規則の反復更新と制度の循環を調べるには、長期
実行も必要になる。

優先順位は明確である。第一に監査先頭値の外部固定と役割・保存点権限の
処理または OS 境界での隔離、第二に固定制度・元の RQGM・全構成の比較と主要機能の除去実験、
第三に正解付き因果グラフによる責任帰属と影響修復の精度測定、第四に
最終原稿検証の被覆率と統治に要する時間・トークン・金額の測定である。

長期的な問いは、研究速度を保ちながら、結託した役割や誤った基準が
安定してしまう状態を避けられるかである。Constitutional ARI-RQGM は、
この問いを測定できる再現可能な基盤を与える。現時点での主張はより
限定的である。モデルで定義した研究役割は共進化できる一方、台帳の
更新権限、憲法改正、来歴、最終的な証拠検証は、その役割の支配外に
置かれている。
